\documentclass[11pt,a4paper]{article}

\usepackage[utf8]{inputenc}
\usepackage[indonesian]{babel}
\usepackage[a4paper,top=2cm,bottom=2cm,left=2cm,right=2cm]{geometry}
\usepackage{amsmath,amsfonts,amssymb}
\usepackage{graphicx}
\usepackage{booktabs}
\usepackage{enumitem}
\usepackage{xcolor}
\usepackage{tcolorbox}
\usepackage{fancyhdr}
\usepackage{titlesec}
\usepackage{hyperref}

% Color Definitions
\definecolor{primaryblue}{RGB}{24, 43, 73}
\definecolor{secondaryblue}{RGB}{41, 128, 185}
\definecolor{accentblue}{RGB}{41, 128, 185}
\definecolor{boxbg}{RGB}{245, 248, 252}
\definecolor{goldaccent}{RGB}{211, 84, 0}

\pagestyle{fancy}
\fancyhf{}
\fancyhead[L]{\small\textbf{\color{primaryblue}Lembar Kerja Mahasiswa \& Problem Set (Minggu 2)}}
\fancyhead[R]{\small\color{gray}Analisis Deret Waktu -- Teknik Elektro}
\fancyfoot[L]{\footnotesize\color{gray}Program Studi Teknik Elektro FT UNIB}
\fancyfoot[R]{\footnotesize\thepage}

\titleformat{\section}{\color{primaryblue}\normalfont\large\bfseries}{\thesection.}{0.5em}{}[\titlerule]
\setlist{nosep}

\begin{document}

% HEADER
\begin{center}
    {\large \textbf{\color{primaryblue}UNIVERSITAS BENGKULU -- FAKULTAS TEKNIK}}\\[0.1cm]
    {\large \textbf{\color{secondaryblue}JURUSAN TEKNIK ELEKTRO}}\\[0.25cm]
    {\LARGE \textbf{\color{primaryblue}LEMBAR KERJA \& PROBLEM SET MINGGU 2}}\\[0.1cm]
    {\normalsize \textbf{Topik: Stasioneritas, Diferensiasi, dan Uji ADF/KPSS}}
\end{center}

\vspace{0.2cm}

\begin{tcolorbox}[colback=boxbg,colframe=primaryblue,arc=2mm,title=\textbf{Identitas Mahasiswa}]
\small
\begin{tabular}{p{2.5cm}p{5.5cm}p{2.5cm}p{4.5cm}}
\textbf{Nama} & : \dotfill & \textbf{NPM} & : \dotfill \\
\textbf{Kelas/Kelompok} & : \dotfill & \textbf{Tanggal} & : \dotfill \\
\end{tabular}
\end{tcolorbox}

\vspace{0.2cm}

\section{Panduan Pengerjaan}
Problem set ini dirancang berjenjang berdasarkan Taksonomi Bloom (Level C1 sampai C6) untuk mengukur capaian \textbf{Sub-CPMK 2}. Kerjakan secara mandiri, rapi, dan sistematis.

\section{Problem Set (Level C1 -- C6)}

\begin{enumerate}[leftmargin=1.8em, label=\textbf{Soal \arabic*.}]
    
    % C1 - Remembering
    \item \textbf{[C1 -- Mengingat / \textit{Remembering}]} \\
    Sebutkan dan jelaskan konsep stasioneritas dalam deret waktu. Mengapa data runtun waktu yang stasioner sangat penting sebelum melakukan pemodelan statistik?
    \vspace{2.2cm}

    % C2 - Understanding
    \item \textbf{[C2 -- Memahami / \textit{Understanding}]} \\
    Jelaskan perbedaan antara pengujian stasioneritas menggunakan Augmented Dickey-Fuller (ADF) Test dan Kwiatkowski-Phillips-Schmidt-Shin (KPSS) Test dalam konteks hipotesis nol (H0) dan hipotesis alternatif (H1).
    \vspace{2.8cm}

    % C3 - Applying
    \item \textbf{[C3 -- Menerapkan / \textit{Applying}]} \\
    Diberikan data beban listrik mingguan $y_t$. Lakukan diferensiasi orde pertama pada data $y_t = [100, 110, 105, 120, 115]$, dan jelaskan apakah hasil diferensiasi tersebut menghilangkan tren positif.
    \vspace{2.8cm}

    \newpage

    % C4 - Analyzing
    \item \textbf{[C4 -- Menganalisis / \textit{Analyzing}]} \\
    Seorang analis data daya menguji stasioneritas sinyal fluktuasi tegangan gardu. Hasil uji ADF memberikan p-value 0.15 dan uji KPSS memberikan p-value 0.02. Analisis arti statistik dari kedua hasil pengujian tersebut dan apa kesimpulan akhir stasioneritas data tersebut!
    \vspace{3.5cm}

    % C5 - Evaluating
    \item \textbf{[C5 -- Mengevaluasi / \textit{Evaluating}]} \\
    Evaluasilah dampak dari melakukan over-differencing (diferensiasi berlebih) pada sebuah data historis daya reaktif yang sebenarnya sudah stasioner pada orde pertama. Apa yang terjadi pada varians data dan bagaimana pengaruhnya terhadap pemodelan?
    \vspace{3.5cm}

    % C6 - Creating
    \item \textbf{[C6 -- Merancang / \textit{Creating}]} \\
    Rancang sebuah prosedur (algoritma langkah demi langkah) sistematis untuk mengecek dan menangani masalah non-stasioneritas pada data harga listrik harian (Rupiah/kWh) yang memiliki tren inflasi sekaligus efek musiman mingguan, hingga data tersebut siap dimodelkan.
    \vspace{4.5cm}

\end{enumerate}

\end{document}
